\documentclass[12pt,a4paper,oneside]{report}

\usepackage[a4paper,margin=1in]{geometry}
\usepackage[T1]{fontenc}
\usepackage[utf8]{inputenc}
\usepackage{setspace}
\usepackage{graphicx}
\usepackage{booktabs}
\usepackage{float}
\usepackage{hyperref}
\usepackage{titlesec}
\usepackage{xcolor}
\usepackage[protrusion=true,expansion=false]{microtype}

\hypersetup{
    colorlinks=true,
    linkcolor=black,
    urlcolor=black,
    citecolor=black
}

\setstretch{1.25}
\setcounter{secnumdepth}{3}
\setcounter{tocdepth}{3}

\titleformat{\chapter}[display]
  {\centering\bfseries}
  {\Large Chapter \thechapter}
  {0.5em}
  {\Large}

\titleformat{name=\chapter,numberless}[display]
  {\centering\bfseries\Large}
  {}
  {0pt}
  {}

\titlespacing*{\chapter}{0pt}{-12pt}{24pt}
\titlespacing*{\section}{0pt}{18pt}{7pt}
\titlespacing*{\subsection}{0pt}{14pt}{5pt}

\newcommand{\ReportTitle}{Distributed Edge Artificial Intelligence}
\newcommand{\DegreeName}{Bachelor of Technology}
\newcommand{\DepartmentName}{Computer Science \& Engineering}
\newcommand{\InstitutionName}{Barak Valley Engineering College}
\newcommand{\InstitutionAddress}{Nirala, Sribhumi -- 788701, Assam}
\newcommand{\GuideName}{Mr. Rana Singha}
\newcommand{\GuideDesignation}{[Supervisor designation]}
\newcommand{\SubmissionDate}{[Month--Year]}
\newcommand{\AcademicSession}{[Academic session]}
\newcommand{\StudentNames}{%
  \begin{enumerate}\setlength{\itemsep}{2pt}\setlength{\topsep}{4pt}
    \item Anirbar Saikia
    \item Avirup Roy
    \item Hisbul Bahar Mazumdar
    \item Ritesh Das
    \item Syed Yashin Hussain
  \end{enumerate}%
}

\newcommand{\frontmatterchapter}[1]{%
  \chapter*{#1}%
  \addcontentsline{toc}{chapter}{#1}%
}

\newcommand{\InstitutionHeader}{%
  \begin{center}
    \includegraphics[width=1.35in]{images/bvec_logo.png}\par
    \vspace{0.15cm}
    {\large\bfseries\color{red} Department of Computer Science \& Engineering\par}
    \vspace{0.12cm}
    {\Large\bfseries\color{blue} BARAK VALLEY ENGINEERING COLLEGE\par}
    \vspace{0.08cm}
    {\small (A Govt.\ of Assam Institution, Approved by AICTE, New Delhi, Govt.\ of India)\par}
    {\small District: Sribhumi, Assam, 788701\par}
  \end{center}
  \vspace{0.2cm}
  \hrule
  \vspace{0.55cm}
}

\newcommand{\CertificateHeading}[1]{%
  \begin{center}
    {\large\bfseries #1\par}
  \end{center}
  \vspace{0.8cm}
}

\begin{document}

\begin{titlepage}
    \newgeometry{top=1in,bottom=0.4in,left=1in,right=1in}
    \centering
    \setstretch{1.1}
    \vspace*{1.3cm}

    {\Large\bfseries \ReportTitle\par}
    \vspace{1.5cm}

    {\large\itshape A Thesis (Major Project) submitted in partial fulfilment\\
    of the requirements\par}
    \vspace{1.2cm}

    {\large\itshape For the Degree of\\[0.15cm]
    \DegreeName\\[0.15cm]
    In\\[0.15cm]
    \normalfont\bfseries Computer Science and Engineering\par}
    \vspace{1.2cm}

    {\large Submitted by\par}
    \vspace{0.6cm}
    \begin{tabular}{l}
      Anirbar Saikia\\
      Avirup Roy\\
      Hisbul Bahar Mazumdar\\
      Ritesh Das\\
      Syed Yashin Hussain
    \end{tabular}
    \vspace{1.0cm}

    {\large Under the Guidance of\\[0.15cm]
    \textbf{\GuideName}\\
    \GuideDesignation\par}
    \includegraphics[width=1.4in]{images/bvec_logo.png}\par
    \vspace{0.25cm}

    {\large\bfseries DEPARTMENT OF COMPUTER SCIENCE \& ENGINEERING\\
    BARAK VALLEY ENGINEERING COLLEGE\\
    \normalsize NIRALA, SRIBHUMI -- 788701, ASSAM\\[0.25cm]
    \large \SubmissionDate\par}
\end{titlepage}
\restoregeometry

\pagenumbering{roman}

\frontmatterchapter{Declaration}

We, the students of \DepartmentName, \InstitutionName, declare that the work
entitled ``\ReportTitle'' has been completed under the guidance of
\GuideName, \GuideDesignation, \DepartmentName, \InstitutionName.

\par\medskip

The report has been prepared without resorting to plagiarism. We have adhered
to the principles of academic honesty and integrity. No falsified or fabricated
data have been presented in this report. Further, the matter embodied in this
project report has not been submitted previously for the award of any degree
or diploma.

\vspace{1cm}

\noindent
Anirbar Saikia\\[0.25cm]
Avirup Roy\\[0.25cm]
Hisbul Bahar Mazumdar\\[0.25cm]
Ritesh Das\\[0.25cm]
Syed Yashin Hussain

\vfill

Date: \rule{5cm}{0.4pt}\hfill Place: \rule{5cm}{0.4pt}

\clearpage
\thispagestyle{plain}
\phantomsection
\addcontentsline{toc}{chapter}{Certificate of Approval}
\InstitutionHeader
\CertificateHeading{CERTIFICATE OF APPROVAL}

\noindent
\textbf{Project Title:} \ReportTitle\\
\textbf{Project Category:} Major Project Report\\
\textbf{Student Names:}

\StudentNames

\noindent
\textbf{Name of the Department:} Computer Science \& Engineering\\
\textbf{Name of the Supervisor:} \GuideName\\
\textbf{Academic Session:} \AcademicSession\\
\textbf{Recommendation:} [Yes/No]

\vfill

\noindent
\textbf{[Name of Head of Department]}\\
Head, Department of Computer Science \& Engineering\\
Barak Valley Engineering College

\clearpage
\thispagestyle{plain}
\phantomsection
\addcontentsline{toc}{chapter}{Certificate from External Examiner}
\InstitutionHeader
\CertificateHeading{CERTIFICATE FROM EXTERNAL EXAMINER}

This is to certify that I have examined the major project report entitled
``\ReportTitle'', submitted by the following students of Bachelor of
Technology in Computer Science and Engineering:

\StudentNames

The report is submitted in partial fulfilment of the academic requirements
for the degree.

\vfill

\noindent External Examiner: \rule{7cm}{0.4pt}\\[0.8cm]
\noindent Signature: \rule{5cm}{0.4pt}\hfill Date: \rule{4cm}{0.4pt}

\clearpage
\thispagestyle{plain}
\phantomsection
\addcontentsline{toc}{chapter}{Certificate from Supervisor}
\InstitutionHeader
\CertificateHeading{CERTIFICATE FROM SUPERVISOR}

This is to certify that the major project/thesis work embodied in this report
entitled ``\ReportTitle'', submitted by:

\StudentNames

of Bachelor of Technology in the Department of Computer Science \& Engineering
is based on their work under my supervision and is prepared for their academic
requirements. It is also certified that this work has not been submitted
elsewhere for any degree or diploma.

\vfill

\noindent
Date: \rule{4cm}{0.4pt}\hfill\textbf{\GuideName}\\
Place: \rule{4cm}{0.4pt}\hfill\GuideDesignation

\clearpage
\frontmatterchapter{Abstract}

\noindent\textbf{Formatting sample only --- replace before submission.}
This sample abstract demonstrates the report's paragraph spacing and text
width. Distributed Edge Artificial Intelligence studies how computation and
learning tasks can be placed across devices near data sources and more
centralized computing resources. The final abstract should briefly state the
project problem, objectives, approach, evaluation method, and conclusions.

\par\medskip
\noindent\textit{Replace this sample with a concise abstract based on the
completed project; do not treat this text as a project result.}

\clearpage
\frontmatterchapter{Acknowledgement}

\noindent\textbf{Formatting sample only --- replace before submission.}
The authors would like to thank their supervisor, Mr. Rana Singha, for
guidance during the project. We also acknowledge the support of the Department
of Computer Science \& Engineering and everyone who contributed advice or
feedback.

\par\medskip
\noindent\textit{Update this acknowledgement to reflect the actual
contributions and permissions before submission.}

\clearpage
\tableofcontents
\clearpage
\listoftables
\clearpage
\listoffigures
\clearpage

\pagenumbering{arabic}

\chapter{Introduction}

\noindent\textbf{Formatting sample only --- replace before submission.}
Artificial intelligence applications increasingly need to process information
close to where it is produced. Distributed Edge Artificial Intelligence
considers how data processing and model execution can be shared between
resource-constrained edge devices and other computing nodes. This chapter
introduces the topic and provides sample section styles for the report.

\section{Background}

An edge-based system may include sensors, user devices, gateways, and
computing infrastructure. The placement of a task can affect response time,
network traffic, privacy, and resource use. These considerations motivate
studying the design choices involved in distributing AI workloads.

\section{Problem Statement}

The project investigates how an AI workload can be organized across edge
resources while accounting for device limitations and application
requirements. The final problem statement should identify the specific
application, constraints, and gap addressed by the project.

\section{Objectives}

The following are sample objective statements; revise them to match the agreed
project scope:
\begin{enumerate}
  \item Study the key concepts and architectures used in distributed edge AI.
  \item Define a suitable system design for the selected application.
  \item Evaluate the design using project-specific metrics and experiments.
\end{enumerate}

\section{Scope and Report Organization}

This sample document uses five chapters to demonstrate the report structure.
The final scope, chapter names, and organization should follow the actual
project and the department's requirements.

\chapter{Literature Survey}

\noindent\textbf{Formatting sample only --- replace before submission.}
This chapter should review research and technical material directly relevant
to the selected project. A well-structured survey groups sources by topic,
compares their approaches, and explains how the review informs the project
design. Add verified citations and a bibliography once the sources have been
selected; no reference entries are fabricated in this formatting sample.

\section{Edge Computing and AI}

Introduce the concepts required to understand the project. Define important
terms consistently and distinguish background information from the
contributions made by the project.

\section{Comparison of Existing Approaches}

Use a consistent basis for comparing related work, such as deployment
setting, processing location, evaluation method, and stated limitations.
Support each comparison with an appropriate source.

When several approaches are discussed, organize them around common themes
rather than presenting a sequence of unrelated summaries. A comparison table
can help readers see which assumptions and evaluation settings differ across
the reviewed work.

\section{Research Gap}

Conclude the survey by identifying the specific question that remains
relevant to this project. The gap should be supported by the reviewed sources
and should lead directly to the objectives and design choices described in
the following chapters.

\chapter{System Design and Methodology}

\noindent\textbf{Formatting sample only --- replace before submission.}
This chapter demonstrates how a figure, numbered sections, and a short
procedural list can be presented. The diagram below is a generic illustrative
workflow and does not describe an implemented system.

\section{Conceptual Workflow}

\begin{figure}[H]
  \centering
  \setlength{\fboxsep}{8pt}
  \resizebox{\textwidth}{!}{%
    \begin{tabular}{c@{\hspace{0.5cm}$\longrightarrow$\hspace{0.5cm}}c@{\hspace{0.5cm}$\longrightarrow$\hspace{0.5cm}}c}
      \fbox{\parbox[c][1.1cm][c]{0.22\textwidth}{\centering Data\\source}} &
      \fbox{\parbox[c][1.1cm][c]{0.22\textwidth}{\centering Edge\\processing}} &
      \fbox{\parbox[c][1.1cm][c]{0.22\textwidth}{\centering Application\\output}}
    \end{tabular}%
  }
  \caption{Illustrative workflow used only to demonstrate figure formatting.}
  \label{fig:sample-workflow}
\end{figure}

Figure~\ref{fig:sample-workflow} is a placeholder example. Replace it with a
diagram that accurately represents the proposed or implemented system.

\section{Sample Method}

An example methodology outline is:
\begin{enumerate}
  \item Define the application requirements and system constraints.
  \item Select the devices, models, and workload distribution strategy.
  \item Implement the design and document the experimental setup.
  \item Evaluate the system using clearly defined metrics.
\end{enumerate}

\section{Evaluation Plan}

The evaluation plan should describe the test environment, input workload,
baseline, and measurement procedure. Define each metric before reporting
results so that readers can interpret the measurements consistently.

\chapter{Results and Discussion}

\noindent\textbf{Formatting sample only --- replace before submission.}
Present the actual observations from the completed experiments in this
chapter. State the test conditions, report units, and explain the limitations
of the evaluation. The entries in Table~\ref{tab:sample-results} are
intentionally marked as placeholders and are not measured results.

\section{Illustrative Results Table}

\begin{table}[H]
  \centering
  \caption{Placeholder results table for demonstrating layout only.}
  \label{tab:sample-results}
  \begin{tabular}{lcc}
    \toprule
    Evaluation case & Metric 1 & Metric 2 \\
    \midrule
    Case A & TBD & TBD \\
    Case B & TBD & TBD \\
    \bottomrule
  \end{tabular}
\end{table}

\section{Discussion}

Replace the placeholder table with measured data and discuss the observed
trends in relation to the project objectives. Avoid drawing conclusions that
are not supported by the evaluation.

\section{Limitations}

Document constraints that may affect interpretation, such as the size of the
test workload, the devices available, or the duration of the experiments.
Explain how these constraints were handled and avoid generalizing beyond the
conditions that were tested.

\chapter{Conclusion and Future Work}

\noindent\textbf{Formatting sample only --- replace before submission.}
The final chapter should summarize what was completed and relate the main
conclusions to the stated objectives. It should distinguish demonstrated
outcomes from proposed future work. This sample contains no project-specific
findings.

Keep the conclusion focused: briefly restate the problem, summarize the
approach, and state only findings established by the report. Avoid introducing
new evidence or claims that have not been discussed in the results.

\section{Future Work}

Possible future-work topics should follow from the actual results and
limitations of the project. Replace this paragraph with relevant next steps
after the evaluation is complete.

\frontmatterchapter{References}

[Add the verified references used in the report. Apply the citation style
required by the department.]

\end{document}
